\documentclass[10pt]{article}
\usepackage[margin=1.8cm, top=1.5cm, bottom=1.5cm]{geometry}
\usepackage{amsmath, amssymb, amsthm}
\usepackage{titlesec}
\usepackage{hyperref}
\usepackage{enumitem}
\usepackage{booktabs}
\usepackage{tabularx}
\usepackage{microtype}
\usepackage{setspace}
\usepackage{graphicx}
\usepackage{fancyvrb}
\usepackage{xcolor}
\usepackage{caption}
\usepackage{subcaption}
\usepackage{float}
\usepackage{array}
\usepackage{multirow}

\hypersetup{colorlinks=true, linkcolor=black, citecolor=black, urlcolor=black}

\titlespacing{\section}{0pt}{6pt}{2pt}
\titlespacing{\subsection}{0pt}{4pt}{1pt}
\titleformat{\section}{\normalfont\bfseries\normalsize}{}{0em}%
  {\thesection\quad}[\vspace{1pt}\hrule\vspace{2pt}]
\titleformat{\subsection}{\normalfont\bfseries\small}{}{0em}{\thesubsection\quad}
\titleformat{\subsubsection}{\normalfont\itshape\small}{}{0em}{\thesubsubsection\quad}

\setlist[itemize]{noitemsep, topsep=2pt, leftmargin=1.2em, parsep=0pt}
\setlist[enumerate]{noitemsep, topsep=2pt, leftmargin=1.5em, parsep=0pt}

\setlength{\parskip}{2pt}
\setlength{\parindent}{0pt}
\captionsetup{font=small, labelfont=bf, skip=4pt}

% ─── convenience macros ──────────────────────────────────────────────────────
\newcommand{\supported}{\textbf{Supported}}
\newcommand{\partialsupport}{\textbf{Partially Supported}}
\newcommand{\notsupported}{\textbf{Not Supported}}
\newcommand{\untestable}{\textit{Untestable}}

% ─── title block ─────────────────────────────────────────────────────────────
\title{\vspace{-1.2cm}
  \textbf{Social Media and Student Well-Being:\\[2pt]
  Second Interim Report}
  \vspace{-0.4cm}}
\author{}
\date{}

% ─────────────────────────────────────────────────────────────────────────────
\begin{document}
\maketitle
\thispagestyle{empty}

% ─── Abstract ────────────────────────────────────────────────────────────────
\begin{abstract}
\noindent
This report examines the relationship between social media usage and student
well-being in a cross-sectional sample of 705 university-aged respondents
(16--25 years). Using a cleaned Kaggle dataset, the eight hypotheses are tested
spanning academic performance, sleep, addiction, platform type, relationship
status, and gender. Exploratory data analysis reveals strong bivariate
associations, formalised through OLS regression, Logit and Probit modelling,
$t$-tests, one-way ANOVA, and a Davidson--MacKinnon
$J$-test. Six of eight hypotheses are supported or partially supported. Daily
social media usage is a strong predictor of both self-reported academic impact
and reduced sleep duration; lower sleep duration is also significantly
associated with reported academic impact; and short-form video platforms
(TikTok, Instagram) are associated with
significantly higher addiction scores. The Logit classifier achieves 91\,\%
accuracy (AUC~$= 0.96$) in predicting self-reported academic impact. Hypotheses
on relationship status and gender are not supported; one hypothesis on
time-of-day usage could not be tested due to data limitations.
\end{abstract}

% ════════════════════════════════════════════════════════════════════════════
\section{Topic and Motivation}
% ════════════════════════════════════════════════════════════════════════════

Social media has become an indispensable part of daily life for students aged
16--25. Platforms such as TikTok, Instagram, and Snapchat are engineered for
sustained engagement, raising legitimate concerns about their impact on academic
performance, sleep quality, and mental health. Despite a growing body of
literature, most existing studies focus on narrow age ranges, single-country
contexts, or self-report measures without accompanying quantitative outcomes.

This project is motivated by lived experience: as students, the
effects of heavy social media use daily is observed---late-night scrolling, shortened
attention spans, and the diffuse anxiety of constant connectivity. The aim is to
translate these observations into rigorous, data-driven analysis. Specifically,
next things are examined:
\begin{itemize}
  \item how much time students spend on social media and which platforms
        dominate;
  \item how usage patterns correlate with academic performance and sleep
        duration;
  \item whether platform type (short-form video vs.\ other) modulates addiction
        risk;
  \item whether demographic factors (gender, relationship status) differentiate
        usage or outcomes.
\end{itemize}

Findings are intended to be useful to students preparing for high-stakes
examinations, school psychologists, university welfare offices, and researchers
working on digital-health interventions.

% ════════════════════════════════════════════════════════════════════════════
\section{Aims, Objectives, and Research Questions}
% ════════════════════════════════════════════════════════════════════════════

\subsection{Central Aim}

To determine \textit{whether, how, and to what degree} social media usage is
associated with student well-being outcomes, using a quantitative
cross-sectional approach.

\subsection{Research Questions}

\begin{enumerate}[label=\textbf{RQ\arabic*.}]
  \item Does spending more time on social media correlate with lower
        self-reported academic performance?
  \item Does greater daily use predict shorter sleep duration?
  \item Is reduced sleep associated with reported academic impact after
        controlling for social media usage?
  \item Are short-form video platforms more strongly associated with addiction
        and poor sleep?
  \item Does relationship status relate to time spent online?
  \item Do students using social media primarily late at night report worse
        outcomes? \textit{(Data limitation: untestable.)}
  \item Do female students report worse well-being effects at equivalent usage
        levels?
  \item Do students with higher addiction scores experience more interpersonal
        conflicts?
\end{enumerate}

\subsection{Hypotheses}

Each hypothesis is stated directionally and corresponds to one research
question.

\medskip
\noindent\textbf{H1.} Students who spend more hours per day on social media
report lower academic performance than those who spend less time.

\medskip
\noindent\textbf{H2.} Higher daily social media usage is negatively correlated
with average sleep duration.

\medskip
\noindent\textbf{H3.} Lower sleep duration is associated with a higher
probability of self-reported academic impact, even after controlling for daily
social media usage and demographics.

\medskip
\noindent\textbf{H4.} Short-form video platforms (TikTok, Instagram) are more
strongly associated with high addiction scores and low grades than long-form or
text-based platforms.

\medskip
\noindent\textbf{H5.} Students in a relationship spend fewer hours on social
media per day than single students.

\medskip
\noindent\textbf{H6.} Students who use social media primarily late at night
report significantly worse sleep and lower grades than those using it earlier
in the day. [\textit{Untestable: no time-of-day variable in the dataset.}]

\medskip
\noindent\textbf{H7.} Female students are more likely to report negative
well-being effects from social media use than male students, even after
controlling for total usage hours.

\medskip
\noindent\textbf{H8.} Students with higher addiction scores report more
conflicts with family or partners related to social media use.

% ════════════════════════════════════════════════════════════════════════════
\section{Literature Review}
% ════════════════════════════════════════════════════════════════════════════

The analytical framework is grounded in three peer-reviewed venues, each
contributing a distinct theoretical lens.

\textbf{Computers in Human Behavior} (Elsevier) provides the primary backbone
for this work on distraction, social comparison, and academic outcomes. Studies
in this journal consistently show that passive social media consumption---
scrolling feeds without direct interaction---is more strongly associated with
negative affect and lower academic engagement than active use.

\textbf{International Journal of Mental Health and Addiction} (Springer) informs
the operationalisation of problematic social media use. The Bergen Social Media
Addiction Scale conceptualises addiction across six dimensions (salience, mood
modification, tolerance, withdrawal, conflict, relapse), mapping closely onto
the \texttt{Addicted\_Score} variable in the dataset.

\textbf{Cyberpsychology, Behavior, and Social Networking} (Mary Ann Liebert)
grounds the platform-specific hypotheses (H4) and this examination of how online
behaviour displaces sleep. Research published here documents the unique
properties of algorithmically curated short-form video feeds---infinite scroll,
variable reward, and high stimulation density per second---believed to make
platforms such as TikTok disproportionately habit-forming relative to earlier
social media formats.

% ════════════════════════════════════════════════════════════════════════════
\section{Data Description}
% ════════════════════════════════════════════════════════════════════════════

\subsection{Source and Structure}

An anonymised cross-sectional dataset sourced from Kaggle is used, titled
\textit{Students Social Media Addiction}. The raw dataset contains 705 student
records across 13 variables. Each row represents one respondent; there are no
repeated measures.

\subsection{Ethical and Geopolitical Exclusions}

All records originating from Russia or Belarus were systematically excluded,
and entries referencing VKontakte were likewise removed. This exclusion resulted
in the removal of a small number of rows. The final cleaned dataset is used for
all analyses reported below.

\subsection{Variables}

Table~\ref{tab:variables} describes all 13 variables, their type, and their
role in the analysis.

\begin{table}[H]
\centering
\caption{Dataset variables and analytical roles.}
\label{tab:variables}
\small
\begin{tabularx}{\linewidth}{@{} l l l X @{}}
\toprule
\textbf{Variable} & \textbf{Type} & \textbf{Scale} & \textbf{Role / Notes} \\
\midrule
Student\_ID & Identifier & --- & Index; excluded from all models \\
Age & Numeric & 16--30 & Demographic control \\
Gender & Categorical & Male/Female & Dummy-coded (\textit{Female}~$= 1$) \\
Academic\_Level & Categorical & High/Under/Grad & Dummy-coded \\
Country & Categorical & 110 countries & Used for filtering; not modelled \\
Most\_Used\_Platform & Categorical & 10 platforms & Grouped into \textit{Short-Form} vs.\ \textit{Other} \\
Avg\_Daily\_Usage\_Hours & Numeric & 0.5--9.9 hrs & Primary predictor across all models \\
Sleep\_Hours\_Per\_Night & Numeric & 4--9 hrs & Outcome (H2) and predictor in Logit B (H3) \\
Mental\_Health\_Score & Numeric & 1--10 & Well-being indicator \\
Addicted\_Score & Numeric & 1--10 & Addiction severity; predictor for H8 \\
Affects\_Academic\_Performance & Binary & Yes/No & Primary outcome (H1, H3); coded 0/1 \\
Conflicts\_Over\_Social\_Media & Numeric & 0--20 & Outcome for H8 \\
Relationship\_Status & Categorical & Single/Partnered/Complicated & Grouping variable for H5 \\
\bottomrule
\end{tabularx}
\end{table}

\subsection{Data Quality}

The dataset is exceptionally clean: zero missing values, zero duplicate rows,
and all numeric variables fall within their stated plausible ranges. No
imputation was required. The \texttt{Student\_ID} column was set as the
DataFrame index and removed from all models.

\subsection{Descriptive Statistics}

\begin{table}[H]
\centering
\caption{Descriptive statistics for numeric variables ($n = 705$ after
exclusions).}
\label{tab:desc}
\small
\begin{tabular}{@{} l r r r r r r @{}}
\toprule
\textbf{Variable} & \textbf{Mean} & \textbf{SD} & \textbf{Min} & \textbf{Q1}
  & \textbf{Median} & \textbf{Max} \\
\midrule
Age & 20.7 & 2.1 & 18 & 19 & 21 & 24 \\
Avg\_Daily\_Usage\_Hours & 4.6 & 2.3 & 0.5 & 2.8 & 4.9 & 9.9 \\
Sleep\_Hours\_Per\_Night & 6.1 & 1.2 & 4.0 & 5.1 & 6.2 & 9.0 \\
Mental\_Health\_Score & 5.4 & 2.5 & 1 & 3 & 5 & 10 \\
Addicted\_Score & 6.4 & 2.2 & 1 & 5 & 7 & 10 \\
Conflicts\_Over\_Social\_Media & 3.1 & 2.8 & 0 & 1 & 3 & 20 \\
\bottomrule
\end{tabular}
\end{table}

Several features of the descriptive statistics warrant attention. Mean daily
usage of 4.6 hours is substantial, with the upper tail reaching 9.9 hours.
Mean sleep of 6.1 hours falls below the 7--9 hour recommendation for this age
group. Addiction scores skew toward higher values (mean 6.4, median 7),
indicating that most respondents self-report moderate-to-high addictive
engagement. Instagram and TikTok are the two most-used platforms, and the
sample spans approximately 110 unique countries with no single country
dominating.

% ════════════════════════════════════════════════════════════════════════════
\section{Exploratory Data Analysis}
% ════════════════════════════════════════════════════════════════════════════

All visualisations were produced in Python using \texttt{pandas},
\texttt{matplotlib}, and \texttt{seaborn}.

\subsection{Univariate Distributions}

Age is roughly uniformly distributed across 18--24, reflecting the
college-aged target population. Daily usage is spread across the full 0.5--10
hour range with a slight right tail. Sleep hours cluster around 5--7, with a
substantial proportion of students sleeping below the clinically recommended
threshold. Addiction scores are markedly right-skewed, with most students
rating themselves 6--9 out of 10. Instagram and TikTok dominate platform
usage, followed by Snapchat and Twitter.

% Figure placeholders --- replace \textit{[Figure X]} with \includegraphics
% once PNG files are available.

\textit{[Figure 1: Three-panel histogram of Age, Avg\_Daily\_Usage\_Hours, and
Sleep\_Hours\_Per\_Night.]}

\textit{[Figure 2: Frequency distributions of Mental\_Health\_Score,
Addicted\_Score, and Conflicts\_Over\_Social\_Media.]}

\textit{[Figure 3: $2\times2$ bar chart grid for Gender, Academic Level,
Most Used Platform, and Relationship Status.]}

\textit{[Figure 4: Top-15 countries by respondent count (horizontal bar
chart).]}

\subsection{Bivariate Analysis}

\subsubsection{H1 --- Usage and Academic Performance}

Students who report academic impact use social media for a median of 5.5 hours
per day, compared to 3.9 hours for those who do not---a difference of 1.6
hours. Density distributions for the two groups are clearly separated, with the
``affected'' group shifted substantially to the right.

\textit{[Figure 5: Boxplot and overlaid density of Avg\_Daily\_Usage\_Hours by
Affects\_Academic\_Performance (Yes/No).]}

\subsubsection{H2 --- Usage and Sleep}

The Pearson correlation between daily usage and sleep is $r = -0.79$, one of
the strongest pairwise relationships in the dataset. Each additional hour of
social media use is associated with meaningfully shorter sleep, as shown by the
downward-sloping OLS regression line.

\textit{[Figure 6: Scatter plot of Avg\_Daily\_Usage\_Hours vs.\
Sleep\_Hours\_Per\_Night with linear fit ($r = -0.79$).]}

\subsubsection{H4 --- Platform Type and Addiction}

TikTok and Instagram users cluster at higher addiction scores, while users of
text-based and long-form platforms (Twitter, YouTube, LinkedIn) display lower
medians.

\textit{[Figure 7: Side-by-side boxplots of Addicted\_Score and
Avg\_Daily\_Usage\_Hours by Most\_Used\_Platform.]}

\subsubsection{H5 --- Relationship Status and Usage}

The three relationship-status groups are virtually indistinguishable: Single
(4.80~hrs), In Relationship (4.70~hrs), Complicated (4.75~hrs). No visual
support for H5 is apparent.

\textit{[Figure 8: Boxplot of Avg\_Daily\_Usage\_Hours by
Relationship\_Status.]}

\subsubsection{H7 --- Gender and Well-being}

Median values across gender groups are nearly identical for daily usage, mental
health score, and sleep duration. The conditional gender effect predicted by H7
requires regression interaction terms, tested in the modelling stage below.

\textit{[Figure 9: Three-panel boxplot of Usage Hours, Mental Health Score, and
Sleep Hours by Gender.]}

\subsubsection{H8 --- Addiction and Conflicts}

The correlation between addiction score and conflicts is $r = 0.93$---the
strongest pairwise relationship in the dataset---with a near-linear positive
association plainly visible in the scatter plot.

\textit{[Figure 10: Scatter plot (with jitter) of Addicted\_Score vs.\
Conflicts\_Over\_Social\_Media with OLS regression line.]}

\subsection{Correlation Structure}

Table~\ref{tab:corr} summarises the key pairwise correlations; the full
annotated heatmap is presented in Figure~11.

\begin{table}[H]
\centering
\caption{Key Pearson correlations between numeric variables.}
\label{tab:corr}
\small
\begin{tabular}{@{} l l r l @{}}
\toprule
\textbf{Variable A} & \textbf{Variable B} & \textbf{$r$}
  & \textbf{Relevance} \\
\midrule
Avg\_Daily\_Usage\_Hours & Sleep\_Hours\_Per\_Night & $-0.79$ & H2 \\
Addicted\_Score & Conflicts\_Over\_Social\_Media & $+0.93$ & H8 \\
Addicted\_Score & Mental\_Health\_Score & $-0.95$ & General well-being \\
Avg\_Daily\_Usage\_Hours & Addicted\_Score & $+0.60$ & General thesis \\
Avg\_Daily\_Usage\_Hours & Mental\_Health\_Score & $-0.55$ & General thesis \\
Age & Avg\_Daily\_Usage\_Hours & $-0.05$ & Negligible \\
\bottomrule
\end{tabular}
\end{table}

\textit{[Figure 11: $6\times6$ annotated Pearson correlation heatmap
(RdBu\_r palette, centred at zero).]}

% ════════════════════════════════════════════════════════════════════════════
\section{Methods}
% ════════════════════════════════════════════════════════════════════════════

Statistical analysis was conducted in Python using \texttt{statsmodels 0.14}
and \texttt{scikit-learn 1.4}. All tests use a significance threshold of
$\alpha = 0.05$ (two-tailed) unless stated otherwise.

\subsection{$t$-tests and ANOVA}

\textbf{Welch's two-sample $t$-test} was used to compare group means for H1,
H4b, and H7. Welch's variant does not assume equal variances, making it more
robust under heteroscedasticity and unequal sample sizes.

\textbf{One-way ANOVA} was applied for H4a (addiction score across all
platforms) and H5 (usage hours across the three relationship-status groups).
ANOVA tests the null hypothesis that all group means are equal.

\subsection{OLS Regression}

Two ordinary least-squares models were estimated:

\begin{enumerate}
  \item \textbf{Model~S (H2).} Sleep duration as the dependent
        variable; predictors: usage hours, age, gender dummy, and academic-level
        dummies.
  \item \textbf{Model~C (H8).} Conflict count as the dependent variable;
        predictors: addiction score, age, and gender dummy.
\end{enumerate}

Both models take the standard form
\begin{equation}
  y_i = \beta_0 + \beta_1 x_{i1} + \cdots + \beta_k x_{ik} + \varepsilon_i,
  \qquad \varepsilon_i \sim \mathcal{N}(0,\sigma^2).
\end{equation}

\subsection{Logit and Probit Models}

Because the primary academic outcome (\texttt{Affects\_Academic\_Performance})
is binary, the probability of a positive response is modeled as
\begin{equation}
  P(y_i = 1 \mid \mathbf{x}_i) = F\!\left(\mathbf{x}_i^\top \boldsymbol{\beta}\right),
\end{equation}
where $F(\cdot)$ is the logistic CDF (Logit) or the standard normal CDF
(Probit). Two Logit specifications are estimated:

\begin{itemize}
  \item \textbf{Logit A (H1).} Predictors: usage hours, age, gender, academic
        level, and a short-form platform dummy.
  \item \textbf{Logit B (H3).} Logit A augmented with
        \texttt{Sleep\_Hours\_Per\_Night} to test whether sleep remains
        significantly associated with reported academic impact after controlling
        for usage and demographics.
\end{itemize}

Marginal effects, evaluated at the mean, are reported alongside coefficients
for substantive interpretation.

\subsection{Model Comparison: $J$-test}

The Davidson--MacKinnon $J$-test is applied to compare the non-nested Logit and
Probit specifications. In $J$-test~1, fitted probabilities from the Probit are
added to the Logit specification; in $J$-test~2, the reverse. Significance of
the added predictions indicates that the alternative model captures information
not present in the primary model.

\subsection{Predictive Evaluation}

The Logit model is evaluated on a held-out test set using an 80/20 stratified
train--test split. Evaluation metrics include a confusion matrix (at threshold
$\tau = 0.5$), precision, recall, $F_1$-score per class, and the ROC curve with
its area under the curve (AUC).

% ════════════════════════════════════════════════════════════════════════════
\section{Results}
% ════════════════════════════════════════════════════════════════════════════

\subsection{H1 --- Usage and Academic Performance}

A Welch $t$-test comparing daily usage between students who report academic
impact ($n \approx 352$, $\bar{x} = 5.9$~hrs, $s = 1.9$) and those who do not
($n \approx 353$, $\bar{x} = 3.3$~hrs, $s = 1.7$) yields $t = 21.5$,
$p \approx 3\times10^{-99}$. The difference is highly significant. Logit A
confirms this at the multivariate level: the marginal effect of one additional
hour of daily usage on the probability of reporting academic impact is
statistically significant after controlling for age, gender, academic level, and
platform type. \textbf{H1 is supported.}

\subsection{H2 --- Usage and Sleep}

OLS Model S explains 66\,\% of the variance in sleep duration ($R^2 = 0.66$,
$F = 271.2$, $p \approx 5\times10^{-161}$). Each additional hour of daily
social media use is associated with a reduction of approximately 0.38~hours in
sleep, controlling for demographics. \textbf{H2 is strongly supported.}

\subsection{H3 --- Sleep and Academic Impact}

In Logit B, \texttt{Sleep\_Hours\_Per\_Night} enters with a negative and
statistically significant coefficient ($\hat\beta = -1.41$, $p < 0.001$).
Holding usage, age, gender, academic level, and platform type constant,
students who sleep less are more likely to report that social media affects
their academic performance. Daily usage remains positive and significant in the
same specification. \textbf{H3 is supported.}

\subsection{H4 --- Short-Form Platforms and Addiction}

One-way ANOVA across all platforms yields $F = 35.7$,
$p \approx 1.5\times10^{-60}$, confirming that addiction scores differ
significantly across platforms. A follow-up Welch $t$-test contrasting
short-form (TikTok, Instagram) against all other platforms yields $t = 9.1$,
$p \approx 1.3\times10^{-18}$, with short-form users reporting meaningfully
higher addiction scores. \textbf{H4 is supported.}

\subsection{H5 --- Relationship Status and Usage}

One-way ANOVA yields $F = 0.41$, $p = 0.66$. Group means are
indistinguishable: Single (4.80~hrs), In Relationship (4.70~hrs), Complicated
(4.75~hrs). \textbf{H5 is not supported.}

\subsection{H6 --- Late-Night Usage}

The dataset contains no time-of-day variable. H6 cannot be evaluated with the
available data. Future research designs should record usage by time window
(e.g., before vs.\ after midnight). \textbf{H6 is untestable.}

\subsection{H7 --- Gender Differences in Well-being}

Welch $t$-tests for \texttt{Mental\_Health\_Score},
\texttt{Avg\_Daily\_Usage\_Hours}, and \texttt{Sleep\_Hours\_Per\_Night} between
male and female respondents all yield $p > 0.05$; median values differ by less
than 0.1 across all three outcomes. Regression models including a gender
$\times$ usage interaction term likewise fail to reach significance.
\textbf{H7 is not supported.}

\subsection{H8 --- Addiction and Conflicts}

OLS Model C achieves $R^2 = 0.874$ ($F \gg 1000$, $p \approx 0$). Each
one-unit increase in addiction score is associated with approximately 0.55
additional conflicts, controlling for age and gender. This is the strongest
explanatory model in the study. \textbf{H8 is strongly supported.}

\subsection{Logit Classifier Performance}

On the stratified 80/20 split, the Logit model achieves accuracy of 91\,\%
and AUC of 0.959, with balanced precision and recall across both classes.
These results confirm that daily usage hours, combined with basic demographic
controls, are strong and generalisable predictors of self-reported academic
impact.

\textit{[Figure 12: Confusion matrix and ROC curve for the held-out test set.]}

\subsection{Logit vs.\ Probit: $J$-test}

Both cross-tests are marginally significant, indicating that neither model
strictly dominates. The Probit attains a slightly lower AIC (386.0 vs.\ 390.1)
and marginally higher pseudo-$R^2$ (0.595 vs.\ 0.591). In substantive terms
the two models are interchangeable; the Logit is retained as the primary model
for the interpretability of its odds ratios.

% ════════════════════════════════════════════════════════════════════════════
\section{Conclusions}
% ════════════════════════════════════════════════════════════════════════════

\subsection{Hypothesis Summary}

\begin{table}[H]
\centering
\caption{Hypotheses, methods, and verdicts.}
\label{tab:summary}
\small
\begin{tabular}{@{} c p{5.4cm} p{3.5cm} c @{}}
\toprule
\textbf{H} & \textbf{Claim} & \textbf{Method(s)} & \textbf{Verdict} \\
\midrule
H1 & More usage $\to$ lower academic performance
   & $t$-test; Logit A & \supported \\
H2 & More usage $\to$ less sleep
   & OLS ($R^2=0.66$) & \supported \\
H3 & Less sleep $\to$ higher reported academic impact
   & Logit B (with sleep and controls) & \supported \\
H4 & Short-form platforms more addictive
   & ANOVA; $t$-test & \supported \\
H5 & Partnered students use SM less
   & ANOVA ($p=0.66$) & \notsupported \\
H6 & Late-night usage $\to$ worse outcomes
   & --- & \untestable \\
H7 & Females report worse effects
   & $t$-tests ($p>0.05$) & \notsupported \\
H8 & Higher addiction $\to$ more conflicts
   & OLS ($R^2=0.87$) & \supported \\
\bottomrule
\end{tabular}
\end{table}

\subsection{Key Findings}

\begin{enumerate}
  \item \textbf{Daily usage is the dominant predictor.} Usage hours is the
        single most important variable across every model estimated. Its
        associations with academic impact, sleep deprivation, and addiction are
        statistically robust and practically meaningful.

  \item \textbf{Sleep is independently associated with academic impact.} Even
        after controlling for daily usage and demographics, lower sleep
        duration is linked to a higher probability of self-reported academic
        impact.

  \item \textbf{Addiction and conflicts are nearly linearly related.}
        $R^2 = 0.87$ for the conflicts model is striking: addiction score
        captures almost all of the explainable variance in interpersonal
        conflict, suggesting that reducing addictive usage patterns may yield
        significant downstream benefits for social relationships.

  \item \textbf{Short-form platforms carry elevated addiction risk.} TikTok and
        Instagram users report addiction scores that are statistically and
        practically higher than users of other platforms. The algorithmic design
        of these platforms---infinite scroll, variable-ratio reinforcement---may
        be the key mechanism.

  \item \textbf{Gender and relationship status do not differentiate outcomes.}
        Neither variable, univariately or in interaction with usage, produces
        significant effects. The risks associated with heavy social media use
        appear to be broadly uniform across gender and relationship status.

  \item \textbf{The Logit classifier generalises well.} AUC of 0.96 and
        accuracy of 91\,\% on a held-out test set confirm that the model
        meaningfully discriminates students who self-report academic harm.
\end{enumerate}

\subsection{Business Relevance and Commercialisation}

Although this project is academic in design, its findings also have practical
market value because they identify measurable behavioural patterns linked to
student well-being and academic risk.

\textbf{Potential clients.}
\begin{itemize}
  \item \textbf{Universities and schools.} Student-support offices, academic
        advising teams, and counselling centres could use this analysis to
        identify high-risk patterns and design prevention programmes.
  \item \textbf{EdTech platforms.} Learning-management systems and study apps
        could integrate risk indicators related to sleep, distraction, and
        problematic social media use.
  \item \textbf{Digital mental-health and wellness companies.} App developers
        focused on student productivity, sleep, or mental health could use the
        model outputs to personalise recommendations and early-warning tools.
  \item \textbf{NGOs and public-sector education bodies.} Organisations working
        on youth mental health or digital literacy may use these results for
        intervention design and evidence-based policy discussion.
\end{itemize}

\textbf{What can be sold.}
\begin{itemize}
  \item \textbf{A reporting product.} A structured institutional report showing
        which usage patterns are associated with academic impact, sleep loss,
        and addiction risk.
  \item \textbf{A dashboard.} A simple analytics dashboard where institutions
        upload student-survey data and receive risk summaries, predicted
        probabilities, and platform-specific comparisons.
  \item \textbf{Consulting and workshops.} Presentation of findings to
        universities, together with recommendations for student well-being
        interventions and digital-hygiene campaigns.
  \item \textbf{A predictive module.} The Logit classifier can be packaged as a
        screening tool for internal use, subject to ethical safeguards and data
        privacy requirements.
\end{itemize}

\textbf{Commercialisation pathway.} The most realistic route is not selling raw
data, but selling \textit{analysis, interpretation, and decision support}. A
small-scale business model could begin with custom reports for universities and
student-support organisations, then expand into a subscription dashboard or a
software module embedded into existing student-success or wellness platforms.
Because the current dataset is cross-sectional and self-reported, the product
should be positioned as a \textit{risk-monitoring and support tool}, not as a
clinical or deterministic decision system.

\subsection{Limitations and Future Directions}

\textbf{Cross-sectional design.} Causality cannot be inferred. Experimental or
longitudinal designs are required to establish whether usage causes harm or
whether students with poorer outcomes gravitate toward heavier use.

\textbf{Self-report measures.} Both addiction scores and academic impact are
subjective. Objective GPA data or device screen-time logs would strengthen
conclusions.

\textbf{Absent variables.} The dataset lacks time-of-day usage data (preventing
H6 testing), objective sleep measures (actigraphy), and clinical mental health
diagnoses.

\textbf{Geographic heterogeneity.} With 110 countries and small per-country
samples, cultural context is inadequately controlled. Country-level fixed
effects would require a larger and more geographically balanced sample.

\textbf{Platform operationalisation.} The binary Short-Form/Other split is a
simplification. Future analyses should apply a finer-grained taxonomy and
account for multi-platform users.

% ════════════════════════════════════════════════════════════════════════════
\appendix
\section*{Appendix: Figure Placement Guide}
\addcontentsline{toc}{section}{Appendix: Figure Placement Guide}
% ════════════════════════════════════════════════════════════════════════════

Table~\ref{tab:figguide} maps each figure placeholder to its source code block
and the file name to be inserted. To insert a figure, replace the corresponding
\textit{[Figure X: \ldots]} placeholder in the \LaTeX{} source with a standard
\verb|\includegraphics| command, for example:

\begin{verbatim}
\includegraphics[width=\linewidth]{figures/fig_h2_sleep.png}
\end{verbatim}

\begin{table}[H]
\centering
\caption{Visualisation placement reference.}
\label{tab:figguide}
\small
\begin{tabularx}{\linewidth}{@{} c l l X @{}}
\toprule
\textbf{Fig.} & \textbf{Section} & \textbf{File to insert}
  & \textbf{Source} \\
\midrule
1 & 5.1 & \texttt{fig\_univar\_continuous.png}
  & Stage 2, Sec.\ 4, histogram block 1 \\
2 & 5.1 & \texttt{fig\_univar\_scores.png}
  & Stage 2, Sec.\ 4, histogram block 2 \\
3 & 5.1 & \texttt{fig\_categorical.png}
  & Stage 2, Sec.\ 4, $2\times2$ subplot \\
4 & 5.1 & \texttt{fig\_countries.png}
  & Stage 2, Sec.\ 4, barh countries \\
5 & 5.2.1 & \texttt{fig\_h1\_academic.png}
  & Stage 2, Sec.\ 5, H1 block \\
6 & 5.2.2 & \texttt{fig\_h2\_sleep.png}
  & Stage 2, Sec.\ 5, H2 scatter \\
7 & 5.2.3 & \texttt{fig\_h4\_platforms.png}
  & Stage 2, Sec.\ 5, H4 boxplots \\
8 & 5.2.4 & \texttt{fig\_h5\_relationship.png}
  & Stage 2, Sec.\ 5, H5 boxplot \\
9 & 5.2.5 & \texttt{fig\_h7\_gender.png}
  & Stage 2, Sec.\ 5, H7 boxplots \\
10 & 5.2.6 & \texttt{fig\_h8\_conflicts.png}
  & Stage 2, Sec.\ 5, H8 scatter \\
11 & 5.3 & \texttt{fig\_heatmap.png}
  & Stage 2, Sec.\ 6, heatmap \\
12 & 7.8 & \texttt{fig\_cm\_roc.png}
  & Stage 3, Sec.\ 6, CM + ROC block \\
\bottomrule
\end{tabularx}
\end{table}

\end{document}
